\uuid{UkJg}
\exo7id{3764}
\titre{exo7 3764}
\auteur{quercia}
\organisation{exo7}
\datecreate{2010-03-11}
\isIndication{false}
\isCorrection{false}
\chapitre{Espace euclidien, espace normé}
\sousChapitre{Endomorphismes auto-adjoints}
\cpge{mpsi-mp2i.matrices-applications-lineaires.a, pcsi.matrices-applications-lineaires.a, ptsi.matrices-applications-lineaires.a, mp-mpi.endomorphismes-euclidien.f, pc.endomorphismes-euclidiens.d, psi.endomorphismes-euclidien.f, pt.isometries-euclidien.e, mp-mpi.endomorphismes-euclidien.b, pc.endomorphismes-euclidiens.b, psi.endomorphismes-euclidien.b, pt.isometries-euclidien.b}
\module{Algèbre}
\niveau{L2}
\difficulte{}

\contenu{
\texte{
Soit ${\cal B} = (\vec e_1,\dots,\vec e_n)$ une base non orthonormée de $E$,
$G$ la matrice de Gram des $\vec e_i$, $f \in \mathcal{L}(E)$ et $M$ sa matrice dans
$\cal B$.
}
\begin{enumerate}
    \item \question{Montrer que $f$ est auto-adjoint si et seulement si ${}^tMG = GM$.}
    \item \question{Montrer que $f$ est orthogonal si et seulement si ${}^tMGM = G$.}
\end{enumerate}
}
